%!TEX root = ../thesis.tex
\chapter{Introduction}
\label{chap:introduction}

\section{The Problem: Baseline Drift in Chromatography}
\label{sec:intro_problem}

Chromatography is among the most widely used analytical techniques in chemistry, pharmacology, environmental science, and biochemistry, enabling the separation and quantification of chemical compounds within complex mixtures. A chromatographic signal---commonly referred to as a chromatogram---records detector response as a function of time or elution volume, producing a series of peaks whose areas are proportional to the concentrations of the corresponding analytes~\cite{skoog2017fundamentals}.

In practice, the raw detector output is corrupted by instrumental and environmental artifacts. The most consequential of these is \emph{baseline drift}: a slowly varying, low-frequency distortion that shifts the signal away from its true zero level. Baseline drift originates from multiple physical sources, including column bleed (thermal degradation of the stationary phase that produces volatile fragments), detector aging and electronic drift, temperature gradients within the chromatographic column, and compositional changes in the mobile phase during gradient elution~\cite{ning2014beads}. The drift is generally smooth and low-frequency relative to the chromatographic peaks, but its precise shape varies unpredictably across instruments, analytes, and operating conditions.

The observed chromatographic signal can be modeled as a three-component additive decomposition:
\begin{equation}
    \yvec = \xvec + \fvec + \wvec,
    \label{eq:signal_model}
\end{equation}
where $\yvec \in \Rspace^N$ is the observed (measured) signal of length $N$, $\xvec \in \Rspace^N$ denotes the \emph{peak component} (the analyte signal of interest), $\fvec \in \Rspace^N$ represents the \emph{baseline} (the low-frequency drift), and $\wvec \in \Rspace^N$ captures additive measurement noise. The central goal of baseline correction is to estimate $\fvec$ from the observed mixture $\yvec$, thereby recovering the peak signal $\xvec$ for downstream quantitative analysis.

Accurate baseline estimation is critical because \emph{peak area integration}---the standard method for determining analyte concentration from chromatograms---is directly affected by the estimated baseline level. An overestimated baseline reduces the computed peak area, leading to systematic underestimation of analyte concentration; conversely, an underestimated baseline inflates peak areas and can produce false positives or concentration overestimates. In regulated domains such as pharmaceutical quality control, food safety testing, and clinical diagnostics, even small systematic errors in baseline estimation can compromise the reliability of analytical results~\cite{ning2014beads, eilers2005baseline}.

% TODO: Uncomment and add figure file when ready.
% Generate with demo.py or demo_chromatogram.py from the Implementations directory.
% \begin{figure}[t]
%     \centering
%     \includegraphics[width=\textwidth]{figures/fig_1_1_chromatogram_drift.pdf}
%     \caption{A representative chromatogram exhibiting baseline drift.
%     \textbf{(a)}~The raw signal $\yvec$ shows visible low-frequency drift
%     obscuring the true peak areas. \textbf{(b)}~After baseline correction,
%     the recovered peak signal $\hat{\xvec}$ enables accurate area integration
%     for quantitative analysis.}
%     \label{fig:chromatogram_drift}
% \end{figure}


\section{Limitations of Existing Approaches}
\label{sec:intro_limitations}

A variety of methods have been proposed for baseline estimation in chromatographic and spectroscopic signals, ranging from polynomial fitting and morphological operations to optimization-based formulations and, more recently, deep learning. Despite decades of development, no existing method simultaneously achieves robust generalization across diverse signal conditions and algorithmic interpretability.

\subsection{Classical Optimization-Based Methods}
\label{sec:intro_classical}

Among classical approaches, BEADS (Baseline Estimation And Denoising using Sparsity), proposed by Ning et al.~\cite{ning2014beads}, is one of the most principled. BEADS formulates baseline correction as a regularized optimization problem that jointly estimates the peak signal and the baseline by exploiting two structural priors: the \emph{sparsity} of the peak signal in the derivative domain (chromatographic peaks occupy a small fraction of the signal) and the \emph{smoothness} of the baseline (drift varies slowly relative to peaks). The formulation integrates a zero-phase Butterworth high-pass filter to enforce frequency-domain separation, first- and second-order total variation penalties to promote sparse derivatives, and an asymmetric penalty controlled by a parameter $r$ that distinguishes positive peaks from negative artifacts.

However, BEADS requires careful specification of six parameters: the sparsity regularization weight $\lambda_0$, the first- and second-order smoothness weights $\lambda_1$ and $\lambda_2$, the asymmetry ratio $r$, the Butterworth filter cutoff frequency $f_c$, and the filter order $d$. In practice, no single parameter configuration generalizes well across instruments, analytes, column conditions, or noise levels. Analysts must tune these parameters for each experimental setup---a tedious and expert-dependent process that undermines the method's practical utility for high-throughput applications.

Other widely used classical methods share this fundamental limitation. Asymmetric Least Squares (AsLS)~\cite{eilers2005baseline} and its adaptive extensions---asymmetrically reweighted Penalized Least Squares (arPLS)~\cite{baek2015arpls} and adaptive iteratively reweighted Penalized Least Squares (airPLS)~\cite{zhang2010airpls}---require specification of smoothing parameters and asymmetry weights. The Statistics-sensitive Non-linear Iterative Peak-clipping (SNIP) algorithm~\cite{morhac2008snip} depends on window size selection. In all cases, optimal parameter values vary across signals, and no automated selection procedure reliably generalizes.

\subsection{Deep Learning Approaches}
\label{sec:intro_dl}

Recent work has applied deep learning to baseline correction, demonstrating the capacity of neural networks to learn the signal-to-baseline mapping directly from training data. Kensert et al.~\cite{kensert2021deep} trained a 1D convolutional autoencoder on 190,000 synthetic chromatograms, achieving effective baseline removal without per-signal parameter tuning. Chen et al.~\cite{chen2022resnet} combined ResNet and U-Net architectures for Raman spectral baseline estimation. More recently, Transformer-based architectures~\cite{zhao20251dtrans} and self-supervised learning frameworks~\cite{hu2024rspssl} have been proposed for spectroscopic baseline correction, further demonstrating the potential of data-driven approaches.

While these methods eliminate the need for per-signal parameter tuning, they introduce a fundamental trade-off: the learned mapping from input signal to estimated baseline is encoded in millions of opaque network parameters, offering no insight into the decomposition process. The resulting models are ``black boxes''---one cannot inspect what regularization strategy the network has learned, verify that it respects physical constraints such as baseline smoothness or peak non-negativity, or diagnose why it fails on a particular signal. This lack of interpretability is a significant limitation in analytical chemistry, where understanding and validating the preprocessing pipeline is often as important as the final result~\cite{gharbi2024unrolled, diras2025}.

\subsection{The Gap}
\label{sec:intro_gap}

The existing landscape thus presents a dichotomy. Classical methods such as BEADS provide interpretable, physically motivated decompositions but require per-signal parameter tuning and cannot generalize across conditions. Deep learning methods achieve generalization through data-driven training but sacrifice the interpretability and structural transparency that domain scientists require. \emph{No existing method bridges this gap}---combining learned adaptivity (eliminating manual tuning) with algorithmic interpretability (retaining the structure and transparency of classical optimization).

Moreover, a fundamental challenge persists even for methods that might bridge this gap: sparsity-based formulations such as BEADS assume that the peak component $\xvec$ is sparse---that chromatographic peaks occupy only a small fraction of the signal length. When peaks are densely packed, as is common in metabolomics and complex mixture analysis, this assumption is violated. The $\ell_1$ penalty over-shrinks dense peaks, and the surplus energy is absorbed into the baseline estimate---a phenomenon we term \emph{baseline leakage}. This is not an implementation artifact but an inherent limitation of sparsity-based signal decomposition, independently confirmed by concurrent work on unrolled sparse-recovery networks~\cite{gharbi2024unrolled} and physics-informed deep learning approaches~\cite{diras2025}. Any learned baseline correction system built on sparse decomposition principles must therefore confront this challenge directly.

Algorithm unrolling~\cite{gregor2010lista, monga2021unrolling} offers a principled framework for closing the interpretability--generalization gap. By mapping each iteration of an iterative optimization algorithm to a neural network layer with learnable parameters, algorithm unrolling creates architectures that inherit the interpretability and inductive biases of the original algorithm while gaining the adaptivity of end-to-end training. This approach has been successfully applied in compressed sensing~\cite{zhang2018istanet}, medical image reconstruction~\cite{yang2018admm}, and other inverse problems, but has seen only limited exploration for chromatographic baseline correction~\cite{gharbi2024unrolled}.


\section{Contributions: LBEADS-NET}
\label{sec:intro_contributions}

This thesis proposes \textbf{LBEADS-NET} (\textbf{L}earned \textbf{B}aseline \textbf{E}stimation \textbf{A}nd \textbf{D}enoising using \textbf{S}parsity \textbf{Net}work), a neural network architecture that bridges the interpretability of classical optimization-based baseline correction with the adaptivity of data-driven learning. The central idea is \emph{algorithm unrolling}: each iteration of the classical BEADS algorithm is mapped to a differentiable neural network layer with learnable parameters, creating a $K$-layer deep network whose architecture directly encodes the optimization structure of BEADS while allowing end-to-end training via backpropagation.

The principal contributions of this thesis are:

\begin{enumerate}
    \item \textbf{Algorithm unrolling of BEADS.} We unroll the BEADS iterative optimization algorithm into $K{=}8$ differentiable layers (BEADSLayers), each with independently learnable regularization parameters $\lambda_0$, $\lambda_1$, $\lambda_2$, and asymmetry ratio $r$---26 learned scalar parameters in total. A log-space parameterization ($\lambda = \exp(\log \lambda)$) guarantees positivity of all regularization weights. The per-layer independence enables \emph{emergent parameter specialization}---early layers learn coarse peak-baseline separation while later layers refine the decomposition---without explicit design at the architectural level.

    \item \textbf{Multi-stage curriculum training with anti-leakage composite loss.} We introduce a three-stage curriculum training strategy to manage an 11-term composite loss function. The loss is designed with two objectives: training accurate signal decomposition and actively suppressing baseline leakage. Key anti-leakage terms include an asymmetric baseline penalty that penalizes over-estimation $9{\times}$ more than under-estimation, an element-wise orthogonality penalty that enforces mutual exclusivity of peak and baseline components, and an envelope constraint that prevents the baseline from exceeding local signal minima. The curriculum progresses from pure reconstruction (Stage~A) through structured separation with anti-leakage losses (Stage~B) to full-loss refinement (Stage~C), stabilizing optimization of an otherwise intractable multi-objective landscape.

    \item \textbf{Systematic diagnosis and mitigation of baseline leakage.} We identify three root causes of baseline leakage in BEADS---sparsity assumption violation in dense-peak regions, low-pass filter bandwidth mismatch, and spatially uniform regularization---and document the systematic mitigation effort across six development iterations. The resulting analysis reveals a fundamental insight: all loss-based mitigations are regularizers on top of a sparsity-biased formulation, and no amount of loss engineering fully overcomes the leakage when the sparsity assumption is violated. This finding generalizes to any unrolled network inheriting $\ell_1$-based sparsity priors, as independently corroborated by concurrent work~\cite{gharbi2024unrolled, diras2025}.

    \item \textbf{Failure mode analysis.} We provide detailed analyses of four additional failure modes: (i)~gradient vanishing in the conjugate gradient solver variant, motivating the ISTA-style reformulation; (ii)~non-learnability of the Butterworth filter cutoff frequency $f_c$ due to the autograd boundary between SciPy and PyTorch; (iii)~training/inference signal length mismatch; and (iv)~softplus approximation artifacts at inference. These findings yield design principles for the algorithm unrolling community---in particular, that the inner solver's gradient properties determine whether an iterative algorithm can be successfully unrolled.

    \item \textbf{Hybrid inference pipeline.} We develop a hybrid inference pipeline that uses the trained LBEADS-NET as a learned initialization for classical BEADS refinement, combined with quality-scored stage selection. This design leverages the network's learned adaptivity for fast approximate decomposition while retaining the convergence guarantees of classical iterative optimization for final refinement.
\end{enumerate}


\section{Thesis Outline}
\label{sec:intro_outline}

The remainder of this thesis is organized as follows. A central thread running through the method design, evaluation, and discussion is the baseline leakage problem---its root causes, the engineering effort to suppress it, and the fundamental limits of that effort.
%
\refchap{chap:background} provides the necessary background, covering chromatographic signal processing, the classical BEADS algorithm in full mathematical detail including its failure modes in dense-peak regions, the algorithm unrolling framework and its origins in LISTA~\cite{gregor2010lista}, and existing neural approaches to baseline correction.
%
Chapter~3 presents the proposed LBEADS-NET architecture, including the mapping from BEADS iterations to differentiable BEADSLayers, the ISTA-style fast variant that resolved gradient vanishing in the conjugate gradient formulation, the multi-stage curriculum training strategy, the eleven-term composite loss function and its anti-leakage design, and the hybrid inference pipeline.
%
Chapter~4 describes the experimental setup: synthetic data generation, the training protocol, evaluation metrics---including leakage-specific measures such as the Baseline Leakage Index and stratified sparse/dense-region evaluation---and baseline comparison methods.
%
Chapter~5 presents quantitative and qualitative results on both synthetic and real chromatographic data, including analysis of learned parameter evolution across layers, baseline leakage metrics, and hybrid pipeline performance.
%
Chapter~6 discusses the principal limitations, with baseline leakage as the centerpiece: the mitigation journey across six development iterations, the fundamental insight that sparsity-biased formulations have inherent limits in dense regimes, and additional failure modes including the conjugate gradient variant's gradient collapse, the non-learnable cutoff frequency, and train/inference gaps.
%
Chapter~7 summarizes the contributions, argues for the generalizability of the approach, and outlines concrete directions for future work motivated by the leakage analysis.
